\documentclass[11pt]{article}
\usepackage[a4paper,margin=20mm]{geometry}
\usepackage{amsmath,amssymb,amsthm,microtype}
\usepackage[colorlinks=true,linkcolor=blue,urlcolor=blue]{hyperref}
\newtheorem{proposition}{Proposition}
\newtheorem{theorem}[proposition]{Theorem}
\newcommand{\id}{\mathbf 1}
\title{Operator Schmidt spaces and circuits for matrix product unitaries}
\author{TNLean contributors}
\date{2 October 2026}
\begin{document}
\maketitle

A possible route to a circuit decomposition of a matrix product unitary is
to express its operator Schmidt factors as unitary operators and then
apply the same construction recursively. This note gives an obstruction
to that requirement. A spin-one reflection gives an obstruction at Schmidt
rank four, and a controlled Pauli construction gives one at Schmidt rank three.
These obstructions concern the proposed decomposition; they do not give
circuit lower bounds. The controlled Pauli example belongs to an explicit
family with linear exact circuits. More generally, classical controls
and a logarithmic quantum target admit polynomial approximate circuits.
The final sections compare positive square-root completions with signed
and square-root-free completions. A final construction establishes a
general quasipolynomial approximation bound, with an error estimate
including the auxiliary output. A stronger exact polynomial theorem,
with exactly erased auxiliaries, is proved in the companion note
\emph{Exact polynomial circuits for matrix product unitaries}, Section~5
(\nolinkurl{docs/audits/2026-10-02_mpu_rank_two_circuits.tex}).
The approximation proof here is independent of that determinant argument.
All circuit statements use arbitrary one- and two-qubit gates.
No claim of priority is intended.

\section{The spin-one reflection}

Let $J_1,J_2,J_3$ be the Cartesian spin-one matrices
\[
 J_1=\begin{pmatrix}0&0&0\\0&0&-i\\0&i&0\end{pmatrix},\quad
 J_2=\begin{pmatrix}0&0&i\\0&0&0\\-i&0&0\end{pmatrix},\quad
 J_3=\begin{pmatrix}0&-i&0\\i&0&0\\0&0&0\end{pmatrix}.
\]
They satisfy $[J_a,J_b]=i\sum_c\varepsilon_{abc}J_c$ and
$\sum_aJ_a^2=2\id$. Write $\sigma_a$ for the Pauli matrices and put
\[
 H=\sum_{a=1}^3J_a\otimes\sigma_a,
 \qquad U=\frac{\id+2H}{3}.
\]

\begin{proposition}\label{prop:reflection}
The operator $U$ is a Hermitian unitary of operator Schmidt rank four.
Its left operator Schmidt space is
\[
 W=\operatorname{span}_{\mathbb C}\{\id,J_1,J_2,J_3\}.
\]
The only unitary matrices in $W$ are $a\id$ with $|a|=1$.
\end{proposition}
\begin{proof}
The Pauli multiplication rule and the spin commutation relations give
\[
 H^2=2\id-H.
\]
For completeness, the off-diagonal part of the Pauli product contributes
\[
 i\sum_{a,b,c}\varepsilon_{abc}J_aJ_b\otimes\sigma_c
 =\frac{i}{2}\sum_{a,b,c}\varepsilon_{abc}[J_a,J_b]\otimes\sigma_c
 =-H.
\]
Thus $H=H^\dagger$ and
$U^2=(\id+4H+4(2\id-H))/9=\id$.

The four left factors $\id,J_1,J_2,J_3$ are linearly independent:
the trace detects the coefficient of $\id$, and the three
independent off-diagonal entries detect the remaining coefficients.
The four right factors $\id,\sigma_1,\sigma_2,\sigma_3$ are also
linearly independent. Their tensor coefficient matrix is diagonal with
four nonzero diagonal entries. This proves the asserted Schmidt rank
and identifies $W$.

Every $X\in W$ has the form $X=a\id+K$, where $K^T=-K$.
An odd-dimensional skew-symmetric matrix has zero determinant:
\[
 \det K=\det K^T=\det(-K)=-\det K.
\]
Choose a nonzero $v\in\ker K$. If $X$ is unitary, then
$Xv=av$ implies $|a|=1$. Since $\operatorname{tr}K=0$,
\[
 3=\operatorname{tr}(X^\dagger X)
   =3|a|^2+\operatorname{tr}(K^\dagger K).
\]
It follows that $K=0$. Conversely, every $a\id$ with $|a|=1$ is unitary.
\end{proof}

\begin{proposition}\label{prop:no-basis}
There is no operator Schmidt decomposition of $U$ with four unitary
left factors. More generally, $W$ is not spanned by its unitary elements.
The same conclusion holds after multiplying $W$ on either side by
fixed unitary matrices.
\end{proposition}
\begin{proof}
The span of the unitary elements of $W$ is exactly
$\mathbb C\id$, whereas $\dim W=4$. In any decomposition
$U=\sum_{a=1}^4A_a\otimes B_a$ with independent right factors,
the left factors span $W$. Multiplication on either side by a unitary
preserves membership in the unitary group and is an invertible linear
map on the matrix space.
\end{proof}

Consequently, a general recursive circuit theorem needs a mechanism
which handles nonunitary factors, or an enlargement of the local
operator spaces. The present example does not exclude a larger linear
combination of unitary product operators: a sufficiently large matrix
algebra is spanned by unitaries. It excludes the more restrictive
proposal that the original Schmidt space always has a unitary basis.
An enlargement must also control the circuit complexity of its new
factors; the mere existence of a unitary dilation of a contraction
does not provide that control.

\begin{proposition}\label{prop:qubit}
There is a three-qubit unitary of Schmidt rank four across the
two-qubit/one-qubit cut whose left Schmidt space contains no unitary
matrix.
\end{proposition}
\begin{proof}
Enlarge the left space from $\mathbb C^3$ to $\mathbb C^4$ and put
$\widehat U=U\oplus\id_2$, where the direct sum refers to the left
space. In tensor notation,
\[
 \widehat U=D\otimes\id_2+
    \frac23\sum_{a=1}^3(J_a\oplus0)\otimes\sigma_a,
 \qquad D=\operatorname{diag}(1/3,1/3,1/3,1).
\]
Its four left and four right factors remain independent. Hence its
left Schmidt space is
\[
 \widehat W=\{aD+(K\oplus0):a\in\mathbb C,\ K^T=-K\}.
\]
If such a matrix were unitary, its last coordinate would give $|a|=1$.
The active three-dimensional block is $a\id_3/3+K$.
As above, $K$ has a nonzero kernel vector, and norm preservation on
that vector gives $|a|/3=1$, a contradiction.
\end{proof}

This enlargement strengthens the obstruction: even a single unitary
Schmidt factor is absent from the left Schmidt space. It also places
the example entirely within qubit systems.

\section{A controlled Pauli unitary of Schmidt rank three}

Let $q_0=e_1$, $q_1=e_2$, $q_2=e_3$, and
$q_3=(3/5)e_1+(4/5)e_2$. The real orthogonal matrix
\[
 O=\id_3-\frac23\mathbf 1\mathbf 1^T,
 \qquad \mathbf 1=(1,1,1)^T,
\]
is a Householder reflection. Put $n_j=q_j$ for $0\leq j<4$ and
$n_{j+4}=Oq_j$. Each $n_j\cdot\sigma$ is a Hermitian unitary. Hence
\[
 V=\sum_{j=0}^7(n_j\cdot\sigma)\otimes|j\rangle\langle j|
\]
is a four-qubit unitary, ordered as one target followed by three controls.
All consecutive cut ranks are at most three. At the first cut, the
three Pauli coefficients are independent, as detected by $q_0,q_1,q_2$.
At the second cut, grouping the first control bit gives coefficient
vectors $(q_j,Oq_j)$, which lie in a three-dimensional subspace.
At the last cut, the right operators lie in the two-dimensional
space of diagonal one-qubit matrices.
Equivalently, an explicit open-boundary MPO has bond dimensions
$3,3,2$. Its target letter is the row of the three Pauli entries;
its first control letter is $\delta_{\mathrm{out},\mathrm{in}}O^a$;
its second is
$\delta_{\mathrm{out},\mathrm{in}}[q_{2b},q_{2b+1}]$;
and its last is $\delta_{\mathrm{out},\mathrm{in}}e_c$, where
$a,b,c\in\{0,1\}$. Their product has target matrix coefficient
$O^a q_{2b+c}\cdot\sigma$.

Its right Schmidt space at the first cut consists of
$\operatorname{diag}(v\cdot n_j)_{j=0}^7$ with $v\in\mathbb C^3$.
If such a matrix were unitary, the first three entries would give
$|v_1|=|v_2|=|v_3|=1$. Write
$s_{ab}=\operatorname{Re}(v_a\overline{v_b})$.
The entry corresponding to $q_3$ gives $s_{12}=0$.
Those corresponding to $Oe_1$ and $Oe_2$ then give
\[
 s_{13}=2s_{23},\qquad 2s_{13}=s_{23}.
\]
Thus all three $s_{ab}$ vanish. These would be three nonzero pairwise
orthogonal vectors in the two-dimensional real plane $\mathbb C$,
which is impossible. The right Schmidt space therefore contains
no unitary matrix. The construction has rational coefficients;
its eight vectors, multiplied by $15$, are
\[
\begin{gathered}
 (15,0,0),\ (0,15,0),\ (0,0,15),\ (9,12,0),\\
 (5,-10,-10),\ (-10,5,-10),\ (-10,-10,5),\ (-5,-2,-14).
\end{gathered}
\]

This example is itself a controlled unitary. It shows that the
requirements of a controlled decomposition and of a unitary basis
in each original Schmidt space are different: the eight orthogonal
control projectors need not lie in the three-dimensional Schmidt space.

\begin{proposition}\label{prop:orthogonal-walk}
Let $n_0\in\mathbb R^3$ be a unit vector and
$O_1,\ldots,O_n\in O(3)$. Put
\[
 n(x)=O_1^{x_1}\cdots O_n^{x_n}n_0,
 \qquad
 U=\sum_{x\in\{0,1\}^n}
   (n(x)\cdot\sigma)\otimes|x\rangle\langle x|.
\]
This unitary has an open-boundary MPO of maximal bond dimension at
most three. It has an exact circuit with $2n$ two-qubit gates,
at most $n+1$ one-qubit gates, and no auxiliary qubits.
With the target at one end of a line, it also has an exact circuit
of size and depth $O(n)$ using adjacent two-qubit gates.
\end{proposition}
\begin{proof}
The MPO is the row of Pauli entries at the target, followed by
control letters $\delta_{\mathrm{out},\mathrm{in}}O_i^{x_i}$,
and the right boundary vector $n_0$.
For each $i$, put $s_i=\det O_i\in\{1,-1\}$ and $R_i=s_iO_i$.
Then $R_i\in SO(3)$ and $O_i=s_iR_i$.
Choose a one-qubit unitary $V_i$ satisfying
\[
 V_i(v\cdot\sigma)V_i^\dagger=(R_iv)\cdot\sigma
 \qquad(v\in\mathbb R^3).
\]
For completeness, if $R_i$ is rotation through angle $\theta$ about
the unit axis $u$, such a lift is
\[
 V_i=\cos(\theta/2)\id-i\sin(\theta/2)(u\cdot\sigma).
\]
The Pauli rule
$(u\cdot\sigma)(v\cdot\sigma)=(u\cdot v)\id+
i(u\times v)\cdot\sigma$ verifies both unitarity and the displayed
conjugation relation, by Rodrigues' rotation formula.

Writing $W(x)=V_1^{x_1}\cdots V_n^{x_n}$ gives
\[
 n(x)\cdot\sigma
   =\left(\prod_i s_i^{x_i}\right)
       W(x)(n_0\cdot\sigma)W(x)^\dagger.
\]
Apply the $n$ controlled gates forming $W(x)^\dagger$, then the
single target gate $n_0\cdot\sigma$, then the $n$ controlled gates
forming $W(x)$. Every $s_i=-1$ contributes an additional $Z$ gate
on its control qubit. These phases are retained, since they vary
with the control word. All gates preserve the control values, and
the construction uses no auxiliary register.

For the line construction, initially place the target to the left of
controls $1,\ldots,n$. Visit the controls in this order, applying
controlled $V_i^\dagger$ and then swapping the target with control $i$.
At the far end apply $n_0\cdot\sigma$. Return in the reverse order,
applying controlled $V_i$ and undoing each swap. The forward sequence
acts on the target as
$(V_n^\dagger)^{x_n}\cdots(V_1^\dagger)^{x_1}=W(x)^\dagger$,
and the return sequence acts as $W(x)$. Every logical qubit returns
to its original position. Apply the determinant phases when their
controls are visited. There are $2n$ additional adjacent swaps,
so both gate count and depth remain linear.
\end{proof}

The preceding rank-three obstruction is in this class. Its first
orthogonal matrix is the Householder matrix $O$ above, its second
has columns
\[
 O_2=[e_3,\ (3/5,4/5,0)^T,\ (-4/5,3/5,0)^T],
\]
and $O_3$ is the rotation sending $e_1$ to $e_2$, $e_2$ to $-e_1$,
and fixing $e_3$. Then
$q_{2b+c}=O_2^bO_3^ce_1$ and
$n_{4a+2b+c}=O^aO_2^bO_3^ce_1$.
It therefore has a linear exact circuit despite the absence of any
unitary in its control-side Schmidt space. The determinant phase of
$O$ gives a $Z$ gate on the first control.
Proposition~\ref{prop:orthogonal-walk} is a mathematical circuit
construction, not a claimed Lean formalization.

\section{A sufficient premise for recursion}

A sufficient premise for the proposed recursion would be the following:
at every balanced cut, both Schmidt spaces have bases of unitary matrices
whose coordinate functionals have operator-norm dual norm at most
$\kappa$, with dimension at most $D$; the same property must hold for
every unitary factor which occurs subsequently. The bound on the
coordinate functionals is essential, since a nearly dependent unitary
basis could otherwise introduce arbitrarily large coefficients.
If $A_i,B_j$ are such bases, their coordinate functionals give
\[
 U=\sum_{i,j}c_{ij}A_i\otimes B_j,
 \qquad |c_{ij}|\leq\kappa^2,
 \qquad \sum_{i,j}|c_{ij}|\leq D^2\kappa^2.
\]
Indeed, each coordinate functional extends to the full matrix algebra
with the same norm and is represented there by a trace-class matrix.
Their tensor product consequently has norm at most $\kappa^2$.
An exact linear-combination construction and amplitude amplification
then use a bounded number of the child circuits, depending only on
$D,\kappa$. Each factor remains in its original Schmidt space and
therefore retains the original internal matrix product ranks.
A balanced recursion with this premise throughout has polynomial
circuit size. This is a conditional criterion: Proposition~\ref{prop:qubit}
shows that its unitary-basis premise fails for general MPUs.

\section{Approximation for a small target register}

The controlled examples suggest a positive statement which does not
require unitary bases in their Schmidt spaces.

\begin{proposition}\label{prop:controlled-approx}
Let $n,t$ be positive integers. Suppose an $(n+t)$-qubit unitary
\[
 U=\sum_{y\in\{0,1\}^n}|y\rangle\langle y|\otimes V(y)
\]
has an open-boundary matrix product representation of maximal bond
dimension $D$, where $V(y)$ acts on $t$ target qubits.
For $0<\varepsilon<1$, there is a circuit using arbitrary one- and
two-qubit gates, with gate and auxiliary-qubit counts polynomial in
$n+t,D,2^t,\log(1/\varepsilon)$, which implements
\[
 \widetilde U=\sum_y|y\rangle\langle y|\otimes\widetilde V(y),
 \qquad \|\widetilde U-U\|\leq\varepsilon.
\]
Every $\widetilde V(y)$ is unitary, and all auxiliary qubits return
exactly to $|0\rangle$.
This is a nonuniform circuit-existence statement.
\end{proposition}
\begin{proof}
Put $N=n+t$ and $q=2^t$. The normalized Choi vector of $U$ has a
right-canonical matrix product representation of bond dimension at
most $D$, obtained by successive singular-value decompositions.
Each letter matrix has operator norm at most one. For fixed
$y$ and fixed target indices $a,b$, its contraction gives
$2^{-N/2}V(y)_{ab}$. Repeated matrix--vector multiplication evaluates
this contraction using $O(ND^2)$ scalar arithmetic operations.

Quantize the canonical tensors and the intermediate arithmetic with
$B$ fractional binary digits. The exact partial contractions have
Euclidean norm at most one. Entrywise tensor errors and arithmetic
rounding therefore give a contraction error at most
$CND^2 2^{-B}$, provided $ND2^{-B}$ is small; the slightly enlarged
matrix norms contribute a factor bounded by a constant.
Allowing $O(\log D)$ additional integer digits also bounds all partial
arithmetic sums. After rescaling by $2^{N/2}$, a choice
\[
 B\geq \left\lceil\frac N2\right\rceil+
       \left\lceil\log_2\frac{CND^2q}{\delta}\right\rceil
\]
gives a computed $q\times q$ matrix $W(y)$ with
$\|W(y)-V(y)\|\leq\delta$, uniformly in $y$.
The bounded canonical tensor entries, to this finite precision, are
constants of the circuit. No lower bound on a nonzero Choi Schmidt
value is used.

We describe finite-precision synthesis on the $q$-dimensional target
space, retaining its dependence on $q$. All constants denoted by $C$
below are numerical constants which can be enlarged finitely many
times. Assume $\delta<1/2$, so $\sigma_{\min}(W)\geq1/2$.
Apply QR elimination, numerically choosing a largest remaining
column entry as a row pivot. Compare the numerical computation with
exact QR using precisely the same row permutations. The residual
column of the exact computation has norm at least $1/2$: its
distance from the preceding column span is bounded below by the
smallest singular value. Thus some entry has magnitude at least
$1/(2\sqrt q)$. If the numerical matrix error is at most
$1/(8\sqrt q)$, the numerically chosen largest entry has exact
magnitude at least $1/(4\sqrt q)$. Subsequent elimination steps
only increase the running pivot magnitude. This comparison does
not require different pivot choices to depend continuously on $W$.

To eliminate an entry $b$ below a pivot $a$, use the two-level unitary
\[
 G(a,b)=\frac1r
   \begin{pmatrix}\overline a&\overline b\\-b&a\end{pmatrix},
 \qquad r=\sqrt{|a|^2+|b|^2}.
\]
It maps $(a,b)^T$ to $(r,0)^T$. The pivot lower bound controls the
division and makes this matrix an $O(\sqrt q)$-Lipschitz function
of its input entries. A bound $Cq$ suffices for a full matrix update.
There are at most $q^2$ elimination steps. If each arithmetic update
and each synthesized rotation has error at most $u$, the numerical
matrix errors obey a bound
\[
 E_{k+1}\leq Cq E_k+Cq u,
 \qquad E_k\leq (Cq)^{q^2+1}u.
\]
Choosing $u$ sufficiently small ensures the pivot lower bound
throughout, by induction. The last pivot is made positive by a
one-basis-state diagonal phase; its magnitude is at least $1/2$.

Here is an explicit finite-precision implementation of the rotations.
Write $\alpha=\arg a$, $\beta=\arg b$, and
$\theta=\arctan(|b|/|a|)$. With the usual qubit rotation convention,
\[
 G(a,b)=R_z(\alpha+\beta)R_y(-2\theta)R_z(\alpha-\beta).
\]
If $|b|$ is below a precision threshold $\tau$, replace it by zero;
this changes $G$ by $O(\sqrt q\,\tau)$ and retains
$G(a,0)=\operatorname{diag}(e^{-i\alpha},e^{i\alpha})$.
Set $\beta=0$ in this branch. Perform the elimination gate even
when $b=0$, so each nonfinal pivot is made positive rather than
leaving a diagonal phase unrecorded.
Otherwise $|b|\geq\tau$, so its phase can be computed with angular
error $\rho$ using
$O(\log(q+1)+\log(1/\tau)+\log(1/\rho))$ digits.
For argument extraction, divide the smaller real component by the
larger and record the quadrant. The denominator is then at least
$|b|/\sqrt2$. Square roots and divisions use fixed-point integer
arithmetic: binary bisection gives a square root to $B$ digits in
$O(B)$ comparisons and multiplications, and long division gives the
quotients. For arctangents on $[-1,1]$, the identity
\[
 \arctan x=2\arctan\frac{x}{1+\sqrt{1+x^2}}
\]
reduces the Taylor-series argument to modulus at most
$\sqrt2-1$, so $O(\log(1/\rho))$ terms suffice. Reciprocal range
reduction handles a larger initial argument. All these operations
take polynomially many Boolean gates in the number of digits.
The same computed $\alpha,\beta$ are reused in their sum and
difference. Accuracy modulo $2\pi$ suffices: a branch change gives
two cancelling signs in the two $R_z$ factors. Error estimates follow
from the displayed matrix entries, including at an argument branch cut.

Exact QR elimination gives $W=QR$ with positive diagonal in the
upper triangular matrix $R$. Put
$\gamma=\|R^\dagger R-\id\|\leq3\delta$.
The diagonal entries of $R$ are at least $1/2$: the diagonal entries
of $R^{-1}$ have modulus at most $\|R^{-1}\|\leq2$.
Also $|R_{ij}|\leq2$. Let $h_k$ bound $|R_{kk}-1|$ and all
$|R_{kj}|$ with $j>k$. The Gram identities give
\[
 |R_{kj}|\leq2\gamma+4q\max_{i<k}h_i,
 \qquad
 |R_{kk}-1|\leq\gamma+2q\max_{i<k}h_i.
\]
For the second inequality use
$|R_{ik}|^2\leq2|R_{ik}|$ and the positive diagonal.
Consequently
\[
 h_k\leq(Cq)^k\gamma,
 \qquad \|R-\id\|\leq q(Cq)^q\gamma,
 \qquad
 \|Q-V(y)\|\leq \delta+3q(Cq)^q\delta.
\]
Choose $\delta\leq\varepsilon/[Cq(Cq)^q]$.
Choose $u\leq\varepsilon/(Cq)^{q^2+C}$ and take
$\tau,\rho\leq u/(Cq)$. Enlarging $C$ also covers the comparison
of the finite-precision rotations with the exact same-pivot QR
factors and their product. If the elimination factors form $L$,
then $Q=L^\dagger$; take the adjoints of the synthesized factors
in reverse order. Each finite-angle factor is exactly unitary,
so this product $\widetilde V(y)$ satisfies
$\|\widetilde V(y)-V(y)\|\leq\varepsilon$.
These estimates are deliberately crude. Their small error tolerances
require only
\[
 B_{\mathrm{tot}}=
 O\!\left(N+\log(ND+1)+q^2\log(q+1)+\log(1/\varepsilon)\right)
\]
binary digits, after increasing the constant to account for phase
extraction and division by quantities as small as $\tau$.
Exponential error constants thus give polynomial precision costs.

For completeness, two-level gates can be compiled with costs
polynomial in $t$ and the number of angle bits. A path between two
target basis words which changes one bit at a time has length at
most $t$. Specified-basis transpositions along this path conjugate
the selected pair to two words differing in one bit; undo the
transpositions after the central gate. There are at most
$2(t-1)$ such transpositions, each a pattern-controlled $X$ gate.
Compute a conjunction of its controls with a chain of Toffoli
gates on clean auxiliary bits, use the conjunction as a single
control, and reverse the chain. This costs $O(t)$ two-qubit gates
and $O(t)$ reusable clean bits. The same construction implements
each binary-angle gate, with one extra angle-bit control. A
conservative bound is $O(t^2B_{\mathrm{tot}})$ gates per two-level
factor. A computed row-pivot label can be implemented by testing
each of its $q$ possible values and applying the corresponding
transposition; this adds at most a polynomial factor in $q$.

To see directly that Toffoli needs only a constant number of
two-qubit gates, choose a unitary $A$ with $A^2=X$. Apply controlled
$A$ from the second control, a CNOT from the first to the second,
controlled $A^\dagger$ from the second, the same CNOT, and controlled
$A$ from the first. On control bits $a,b$, the target exponent is
$b-(a\mathbin{\oplus}b)+a=2ab$. Thus this five-gate sequence is
exactly Toffoli, by linearity. One may take
$A=P_++iP_-$ for the spectral projections of $X$.

The last diagonal phase in QR is implemented as a phase on its
target basis state, using the same conjunction construction.
The phase of a two-level $U(2)$ gate can likewise be applied to
the conjunction bit. All these phases are retained;
in particular, a $y$-dependent overall phase of $V(y)$ is preserved.
There are $O(q^2)$ factors. Evaluating all $q^2$ target matrix
entries uses $O(q^2ND^2)$ scalar operations, and QR uses
$O(q^3)$ scalar updates. Binary arithmetic, the rotation
construction, and reversible storage therefore all have polynomial
cost in $N,D,q,\log(1/\varepsilon)$.

All tensor evaluations, pivot choices, and angle calculations depend
only on $y$. Implement their Boolean circuits reversibly, retain their
intermediate bits, apply the controlled target gates, and reverse the
arithmetic circuit. The target gates leave $y$ and the arithmetic
bits unchanged. Thus every auxiliary bit is erased exactly, even on
a superposition of controls. Finally,
\[
 \|\widetilde U-U\|
   =\max_y\|\widetilde V(y)-V(y)\|\leq\varepsilon.
\]
There is no factor proportional to the number of control words.
\end{proof}

In particular, for any fixed $c>0$, the same conclusion has
polynomial cost in $N,D,\log(1/\varepsilon)$ whenever
$t\leq c\log_2 N$. The degree of this polynomial may depend on $c$.

Proposition~\ref{prop:controlled-approx} is a mathematical circuit
construction; it is not asserted to have been formalized in Lean.
It does not supply a polynomial-time numerical conversion of an
arbitrary supplied tensor presentation to its canonical form.

\section{Positive defects and signed completions}

A positive square root need not preserve the matrix-product structure
of an operator Schmidt factor. The following example occurs within
bond-two MPUs and also supplies a completion for which the bond remains
two.

Fix $m\geq1$, put $q=2^m$ and $h=\pi/(2q)$, and consider two consecutive
blocks of $m$ control qubits followed by one target qubit. Write their
computational basis labels as
\[
 i=\sum_{k=0}^{m-1}2^k y_k,\qquad
 j=\sum_{k=0}^{m-1}2^k z_k,\qquad 0\leq i,j<q.
\]
Define a diagonal self-adjoint operator on the controls by
$\Theta|i,j\rangle=h(i-j)|i,j\rangle$, and put
\[
 E=\cos\Theta,\qquad F=\sin\Theta,\qquad
 U=e^{i\Theta\otimes X}=E\otimes\id+iF\otimes X,
\]
where $X$ is Pauli $X$. This unitary is implemented exactly by $2m$
two-qubit gates: conditional rotations $e^{ih2^kX}$ for $y_k=1$ and
$e^{-ih2^kX}$ for $z_k=1$. If $P_\pm$ are the spectral projections of
$X$, then
\[
 U=e^{i\Theta}\otimes P_+ + e^{-i\Theta}\otimes P_-.
\]
The two exponentials are products of one-qubit diagonal unitaries.
Hence $U$ has an open-boundary MPO representation of bond dimension
two in the stated site order. The angles may depend on the site and
on $m$; this is a fixed-bond example in the nonuniform class.

The factors $\id,X$ are Hilbert--Schmidt orthogonal, as are $E,F$:
the identity $\operatorname{Tr}(EF)=0$ follows by exchanging $i$ and
$j$. Both $E$ and $F$ are nonzero. Thus the displayed expansion is an
operator Schmidt decomposition across the controls/target cut after
positive norm rescaling. Moreover,
\[
 E^\dagger E+F^\dagger F
 =EE^\dagger+FF^\dagger=\id.
\]
In particular, $E$ is an actual contraction factor of this MPU.

Its positive defect is
\[
 D_+=\sqrt{\id-E^\dagger E}=|\sin\Theta|.
\]
Across the cut between the two control blocks, its operator Schmidt
rank is the rank of
\[
 K_{ij}=\sin\bigl(h|i-j|\bigr),\qquad 0\leq i,j<q.
\]
This matrix is invertible. Indeed, if $Kv=0$, then for every interior
index $1\leq a\leq q-2$ the sine recurrence gives
\[
 (Kv)_{a+1}-2\cos(h)(Kv)_a+(Kv)_{a-1}
       =2\sin(h)v_a.
\]
Thus all interior components vanish. The first and last rows then give
$\sin((q-1)h)v_{q-1}=\sin((q-1)h)v_0=0$, so $v=0$.
For $q=2$ these endpoint equations already suffice. Consequently
$D_+$ has operator Schmidt rank $q=2^m$ at this cut. By contrast,
\[
 E_{ij}=\cos(hi)\cos(hj)+\sin(hi)\sin(hj),\qquad
 F_{ij}=\sin(hi)\cos(hj)-\cos(hi)\sin(hj)
\]
have rank two.

The positive-defect completion
\[
 H_+=\begin{pmatrix}E&D_+\\D_+&-E\end{pmatrix}
\]
is unitary. Place its one-qubit flag after the controls. Taking the
flag matrix entry with output $0$ and input $1$ extracts $D_+$; this
operation acts on one side of the central physical cut and cannot
increase operator Schmidt rank. Thus $H_+$ has bond dimension at least
$2^m$.

Nevertheless the signed completion
\[
 H_s=\begin{pmatrix}E&-F\\F&E\end{pmatrix}
       =e^{-i\Theta\otimes Y}
\]
is unitary and has bond dimension two. Here $Y$ is Pauli $Y$ on the
flag. It has the same $2m$-gate construction, with the target generator
replaced by $-Y$, and implements
\[
 \psi\otimes|0\rangle\longmapsto
 E\psi\otimes|0\rangle+F\psi\otimes|1\rangle.
\]
The positive square root discarded a sign which was compatible with
the original matrix-product representation. The example therefore
rules out a general bounded-bond assertion for positive-defect
completions; it does not rule out other completions or efficient
circuits.

\section{A square-root-free completion of a Kraus isometry}

Let $A_1,\ldots,A_r$ act on a finite-dimensional Hilbert space
$\mathcal H$ and satisfy
\[
 \sum_{a=1}^r A_a^\dagger A_a=\id.
\]
In $\mathcal H\otimes\mathbb C^{r+1}$ define orthogonal isometries
\[
 J\psi=\psi\otimes|0\rangle,\qquad
 V\psi=\sum_{a=1}^r A_a\psi\otimes|a\rangle.
\]
Then
\[
 S=\id-(V-J)(V-J)^\dagger
   =JV^\dagger+VJ^\dagger+\id-JJ^\dagger-VV^\dagger
\]
is a Hermitian unitary, with $SJ=V$ and $SV=J$. Indeed, if $W=V-J$,
then $W^\dagger W=2\id$, so $(WW^\dagger)^2=2WW^\dagger$ and
$S^2=\id$. The same identities give its action on $J$ and $V$.
On the orthogonal complement of their ranges, $S$ acts as the identity.

Its blocks in the flag basis are particularly simple:
\[
 S_{00}=0,\qquad S_{0a}=A_a^\dagger,\qquad S_{a0}=A_a,
 \qquad S_{ab}=\delta_{ab}\id-A_aA_b^\dagger
 \quad(1\leq a,b\leq r).
\]
Thus no positive square root occurs. If each $A_a$ has operator Schmidt
rank at most $D$ at every internal cut of $\mathcal H$, and the flag
is placed at an endpoint, subadditivity and multiplicativity of these
ranks give the conservative bound
\[
 1+2rD+r^2D^2
\]
for $S$ at every such cut. The flag may be encoded in
$\lceil\log_2(r+1)\rceil$ qubits, extending $S$ by the identity on
unused flag states. This concerns a unitary completion of an isometry;
the output flag of $V$ is part of that isometry.

For an operator Schmidt expansion
$U=\sum_a A_a\otimes B_a$ of a unitary, normalized by
$\operatorname{Tr}(B_a^\dagger B_b)=\dim(\mathcal H_R)\delta_{ab}$,
the normalized partial traces of $U^\dagger U$ and $UU^\dagger$ give
\[
 \sum_a A_a^\dagger A_a=\sum_a A_aA_a^\dagger=\id.
\]
The resulting $V$ preserves the original internal bond bound $D$:
it is obtained by applying a linear map $B_a\mapsto|a\rangle$ on
the opposite side of every internal cut. Consequently its completion
$S$ has the sharper bound $1+2D+D^2=(D+1)^2$ at those cuts.

This construction supplies a controlled bond bound for one completion,
but does not establish a polynomial circuit bound under repetition.
The displayed bound is quadratic in the current bond dimension;
repeating it over a logarithmic number of subdivisions can produce
exponential growth. The Kraus identities alone do not supply a
bounded-dimensional algebra of physical Schmidt factors. In the
preceding example, $E+iF=e^{i\Theta}$ has $2q-1$ distinct eigenvalues,
so the commutative $C^*$-algebra generated by $E,F$ has dimension
$2q-1$, although $r=D=2$. Its spectral projections are polynomials
in $e^{i\Theta}$. A general recursive argument therefore still needs
an additional invariant or a construction which preserves the
appropriate cancellations.


\section{A general quasipolynomial approximation}

The following statement applies to arbitrary fixed-bond MPUs. Its error
includes the auxiliary output: the circuit is close to the desired
isometry on clean auxiliary inputs. It does not assert exact erasure of
those outputs or polynomial circuit size. The construction is nonuniform
and uses elementary amplitude amplification and Fourier series.

\begin{theorem}\label{thm:mpu_general_quasipolynomial_approximation}
Let $N\geq1$, and let $U$ be an $N$-qubit unitary whose operator Schmidt rank across every
consecutive cut is at most $D$. For $0<\varepsilon<1$, there are auxiliary
qubits, an isometry $J\psi=\psi\otimes|0\cdots0\rangle$, and a circuit
unitary $W$ such that
\[
 \|WJ-JU\|\leq\varepsilon.
\]
For fixed $D$, the gate and auxiliary-qubit counts are at most
\[
 \exp\!\left\{C_D\log(N+1)
      \log\bigl(\log(N+1)+\log(1/\varepsilon)+2\bigr)\right\}.
\]
The constant $C_D$ depends only on $D$. The theorem does not assert an
$N$-qubit circuit approximating $U$ with exactly erased auxiliary outputs,
or a bound on $\|W-U\otimes\id\|$ on arbitrary auxiliary inputs.
\end{theorem}

\begin{proof}

\subsection{Contractions at a cut}

Let $A$ be a contraction on a finite tensor product, of operator Schmidt
rank $r$ at a chosen cut. Its left Schmidt space $E$, with the operator
norm, has an Auerbach basis $A_1,\ldots,A_r$: each $\|A_i\|=1$ and its
coordinate functional $\phi_i$ has norm one. One elementary construction
maximizes the absolute determinant of a basis among tuples in the unit
ball. Replacing one basis vector by any vector in that ball proves the
coordinate-functional bound.

Extend $\phi_i$ to the full left matrix algebra without increasing its
norm. Represent it by a trace-class matrix of trace norm one. A singular
value decomposition expresses the resulting slice as a sum of left
matrix-element compressions with total absolute weight one. Consequently
\[
 A=\sum_{i=1}^r A_i\otimes B_i,
 \qquad B_i=(\phi_i\otimes\id)(A),
 \qquad \|A_i\|,\|B_i\|\leq1.
\]
All consecutive internal operator Schmidt ranks of $A_i$ and $B_i$ are
no larger than those of $A$. Indeed, each is a linear slice of $A$ on
the complementary factor. Thus this decomposition may be repeated
on a balanced tree without increasing $D$.

\subsection{An elementary linear amplification polynomial}

Fix once and for all $s=2^{20}$. For $r\geq1$, put
\[
 h=\frac8{rs^2},\qquad K=rs,\qquad H=Kh=\frac8s.
\]
Define an odd, $2\pi$-periodic function
\[
 f_0(\theta)=
 \begin{cases}
 K\sin\theta,&\operatorname{dist}(\theta,\pi\mathbb Z)\leq h,\\
 0,&\text{otherwise}.
 \end{cases}
\]
It is also $\pi$-antiperiodic. It has modulus at most $H$ and support
of total length $4h$ on a period.

For each $j\geq1$, convolve with the uniform probability measure on
$[-a_j,a_j]$, where
\[
 a_j=\frac{h}{20j^2}.
\]
The limiting probability measure $\mu$ is supported on
$[-\sum_j a_j,\sum_j a_j]\subset[-h/10,h/10]$, and its Fourier
coefficients are
\[
 \widehat\mu(n)=\prod_{j\geq1}\operatorname{sinc}(na_j),
 \qquad \operatorname{sinc}(x)=\frac{\sin x}{x}.
\]
The convolution limit can equivalently be taken uniformly after the
first convolution: that first convolution is Lipschitz, and the sum
of all remaining shift widths converges. No distributional limit
is needed for the resulting function. Moreover
\[
 \nu_1:=\widehat\mu(1)\geq
 1-\frac16\sum_j a_j^2\geq\frac12.
\]
Here $\operatorname{sinc}(a_j)\geq1-a_j^2/6$, and the product of
numbers $1-u_j$ with $0\leq u_j\leq1$ is at least $1-\sum_j u_j$.
Define $f=(f_0*\mu)/\nu_1$. Whenever $|\theta|\leq h/4$, all shifts
remain in the central sine interval, and evenness of $\mu$ gives
\[
 f(\theta)=K\sin\theta.
\]

If $j\leq\lfloor\sqrt{nh/40}\rfloor$, then $na_j\geq2$, and
$|\operatorname{sinc}(na_j)|\leq1/2$. The other factors have modulus
at most one. Hence, for $n\geq1$,
\[
 |\widehat\mu(n)|\leq
 2\exp\!\left(-\frac{\sqrt{nh}}{16}\right).
\]
Oddness and antiperiodicity imply that $f$ has only odd sine coefficients:
\[
 f(\theta)=\sum_{\substack{n\geq1\\n\text{ odd}}}b_n\sin(n\theta).
\]
The support and size estimates for $f_0$, followed by convolution and
division by $\nu_1$, give
\[
 |b_n|\leq6Hh\exp\!\left(-\frac{\sqrt{nh}}{16}\right).
\]
These coefficients are absolutely summable. Using
$\int_0^\infty e^{-\sqrt{xh}/16}\,dx=512/h$ yields
\[
 \sum_n|b_n|\leq6H(h+512)\leq3078H<\frac14.
\]
The same integral estimate for the tail gives
\[
 \sum_{n>M}|b_n|
 \leq6144H\exp\!\left(-\frac{\sqrt{Mh}}{32}\right).
\]
In particular, truncation at
$M=O(r(1+\log(1/\delta))^2)$ gives uniform error at most $\delta$
for every $0<\delta<1$.
All constants here are numerical, since $s$ is fixed.

Set
\[
 p_M(x)=\sum_{\substack{1\leq n\leq M\\n\text{ odd}}}
             b_n\sin(n\arcsin x).
\]
For odd $n=2k+1$, $\sin(n\arcsin x)=(-1)^kT_n(x)$ is a polynomial.
The coefficient bound gives $|p_M(x)|\leq1/4$ on $[-1,1]$.
For $|x|\leq1/(rs^2)$, we have
$|\arcsin x|\leq2/(rs^2)=h/4$, and therefore
\[
 |p_M(x)-rsx|\leq\delta.
\]

\subsection{Implementation and perturbation bounds}

Let $P$ be the projection onto zero auxiliary input and output, and let
$Q$ be a circuit unitary with compressed block $Z=PQP$, identified with
a matrix on the data space. Put $R=2P-\id$ and
$G=-QRQ^\dagger R$. A two-dimensional calculation for each singular
value $x$ of $Z$ gives the compressed block of $G^kQ$ as the
singular-value polynomial
\[
 \sin((2k+1)\arcsin x).
\]
For example, the first step gives $3Z-4ZZ^\dagger Z$; the full identity
follows by the plane rotation on the input and output singular vectors.
Thus a linear-combination circuit with weights $b_{2k+1}$ implements
the compressed block $p_M^{\rm sv}(Z)$.

There is no additional normalization: prepare the label register with
amplitudes $\sqrt{|b_n|}$, apply the corresponding controlled $G^kQ$
and its coefficient sign, and undo the preparation. Fill the remaining
probability weight with a branch that flips one additional auxiliary
qubit and hence has zero compressed block. Since $\sum|b_n|\leq1$,
the label preparation is normalized. All real coefficients and label
preparations are nonuniform constants of the circuit.

Here is a useful perturbation estimate that does not require closeness
between unitary dilations. For Hermitian contractions $X,Y$, the
Chebyshev recurrence gives
\[
 \|T_n(X)-T_n(Y)\|\leq n^2\|X-Y\|.
\]
Indeed, writing $E_n=T_n(X)-T_n(Y)$ gives
$E_n=2XE_{n-1}-E_{n-2}+2(X-Y)T_{n-1}(Y)$.
Solving this recurrence using the second-kind polynomials $U_j$,
and using $\|U_j(X)\|\leq j+1$, proves the stated bound.
Apply it to the Hermitian dilation
\[
 \mathcal H(Z)=\begin{pmatrix}0&Z\\Z^\dagger&0\end{pmatrix}.
\]
The off-diagonal block of an odd polynomial in $\mathcal H(Z)$ is
its singular-value polynomial on $Z$. As
$\|\mathcal H(Z)-\mathcal H(Z')\|=\|Z-Z'\|$, it follows that
\[
 \|p_M^{\rm sv}(Z)-p_M^{\rm sv}(Z')\|
 \leq M^2\|Z-Z'\|.
\]

For the gate counts, the controlled powers can share a single loop.
Apply $Q$ once, and for $j=1,\ldots,\lfloor(M-1)/2\rfloor$ apply
$G$ conditioned on the label value $k\geq j$. Each branch then
receives precisely $G^kQ$. Compute and erase each comparison
reversibly using $O(\log M)$ Boolean operations. Give the dummy
branch $k=0$ and flip its separate auxiliary marker; this still has
zero compressed block. Thus there are $O(M)$ calls to
$Q,Q^\dagger$ and the clean projection reflection. The minus sign
in $G$ is a phase on the comparison bit and is retained.
Arbitrary label preparation has polynomial cost in $M$.
A reflection about zero on $a$ auxiliary bits uses
$O(a)$ clean auxiliary bits and $O(a)$ two-qubit gates: compute the
conjunction by a Toffoli chain, apply its phase, and reverse the chain.
Specified label controls are implemented by the same construction.
Controlled versions of one- and two-qubit gates have bounded two-qubit
cost. This last claim also follows directly from QR in dimension eight
and the basis-word path construction for two-level gates. A Toffoli
has a five-two-qubit-gate decomposition using controlled $\sqrt X$
and two CNOT gates. Thus all gate and auxiliary costs here are
polynomial in $M$ and in those of $Q$.

In a recursively controlled circuit there is only one master control
at each invocation. Combine that control with the current label
condition in a fresh guard bit, invoke the child with this one guard
as its master control, and erase the guard. The child does not change
the ancestor control or the current label. This prevents the number
of accumulated controls from growing with the recursion depth.
All conjunction, comparison, and gate-decomposition scratch bits
return exactly to zero for arbitrary states of the data and logical
auxiliary registers. Thus each described $Q$ is a genuine unitary
on those logical registers; its implementation scratch may be
omitted from the mathematical reflection $P$. During repetitions
it remains exactly clean. Reflections inspect logical flags, using
their own clean scratch to compute and erase the conjunction.

\subsection{Balanced recursion}

We recursively approximate the compressed block $A/s$ for every
contraction $A$ of internal cut ranks at most $D$. At a one-qubit leaf,
the standard $4\times4$ unitary dilation of $A/s$ is one arbitrary
two-qubit gate; no recursive square-root construction is used.

At an internal node, use the Auerbach decomposition above and child
circuits whose compressed blocks approximate $A_i/s,B_i/s$ to error
$\eta'$. A uniform linear combination of their products has a
compressed block $Z$ satisfying
\[
 \left\|Z-\frac{A}{rs^2}\right\|\leq2\eta'.
\]
Indeed, actual child blocks are contractions, and the two tensor-product
error terms have norms at most $\eta'$ and $\eta'/s$, respectively.
Choose the Fourier truncation with error $\eta/2$ and
$M\leq C_D(1+\log(1/\eta))^2$. Its singular-value polynomial applied
to the ideal block approximates $A/s$ to error $\eta/2$.
The perturbation estimate shows that taking
\[
 \eta'=\frac{\eta}{12M^2}
\]
keeps the actual compressed-block error below $\eta$.

Let $L=1+\log(1/\eta)$. Along a path down the tree the precision
requirements obey
\[
 L_{j+1}\leq L_j+C_D+4\log(L_j+1).
\]
Consequently $L_j\leq C_D(L_0+j+1)^2$, by induction after enlarging
the constant. More precisely, let $F$ be the sum of the gate and
auxiliary-qubit counts. The controlled constructions above give
\[
 F_{\rm node}\leq C_D M^c\sum_{\rm children}F_{\rm child}
                  +C_D M^c
\]
for a fixed numerical exponent $c$. The logical-flag reflection cost
is included, since its number of flags is bounded by the sum of
child auxiliary counts and the polynomial label-register overhead.
There is no additional factor $N$ at each recursion level.
Controlled sums over the $r\leq D$ choices change only $C_D$.
With $k=\lceil\log_2N\rceil$, this gives
\[
 T(N,\eta),\ a(N,\eta)
 \leq \exp\!\left\{C_D (k+1)\log(L_0+k+2)\right\}.
\]
Zero Schmidt rank means the zero operator and is handled by an auxiliary
flip. Unequal child sizes are split as evenly as possible.

\subsection{The final unitary and the auxiliary output}

For unitary $U$, the construction gives a unitary circuit with
compressed block $Z$ satisfying $\|Z-U/s\|\leq\eta$.
Choose $\ell=\lceil\pi s/4\rceil$, put
$\theta=\pi/(4\ell+2)$, and scale this block by
$\kappa=s\sin\theta\leq1$ using one more auxiliary rotation.
The ideal scaled block is $\sin\theta\,U$. Apply $\ell$ Grover
steps. Its ideal compressed block is exactly $U$, since
$\sin((2\ell+1)\theta)=1$.
The Chebyshev perturbation bound gives an actual block $B$ with
\[
 \|B-U\|\leq(2\ell+1)^2\eta.
\]
All constants in this final step are numerical, independent of $N,D$
and $\varepsilon$.

For the full final circuit $W$ and its clean-input embedding $J$,
unitarity gives
\[
 (WJ-JU)^\dagger(WJ-JU)
   =2\id-B^\dagger U-U^\dagger B.
\]
It follows that
\[
 \|WJ-JU\|\leq\sqrt{2\|B-U\|}.
\]
Take $\eta=\varepsilon^2/[2(2\ell+1)^2]$ to conclude
Theorem~\ref{thm:mpu_general_quasipolynomial_approximation}. The same estimate bounds the nonzero-auxiliary
component of the output. That component need not vanish exactly.

\end{proof}

The kernel-checked identities and the scalar-plus-skew-symmetric
unitarity obstruction are developed in
\nolinkurl{TNLean/MPS/MPU/SpinOneSchmidtObstruction.lean}.
The Pauli witness, its three-term tensor expansion, and the dimension-three
coefficient space containing no unitary are proved in
\nolinkurl{TNLean/MPS/MPU/PauliControlObstruction.lean}.
The orthogonal-isometry exchange and its two actions are proved in
\nolinkurl{TNLean/MPS/MPU/OrthogonalIsometrySwap.lean}.
The physical cut-rank profile and its orthogonal-walk realization are also
checked with exact arithmetic in
\nolinkurl{scripts/check_mpu_pauli_control.py}.
The operator Schmidt rank calculations and the circuit constructions remain
mathematical proofs rather than claimed Lean formalizations.
\end{document}
